% Style: code/trees/cartesian.cpp, vector<int> st and stack indices.
\subsection{单调栈：两侧最近更小或更大元素}
\label{ds:monotonic-stack}
\AlgorithmInfo
    {求每个元素左、右两侧最近的严格更小或严格更大元素的位置。}
    {时间 $O(n)$，空间 $O(n)$。}
    {输入为带 a[0] 占位的 1 下标数组；结果保存下标，不是元素值。}
\subsubsection*{原理}
从左到右扫描，栈中保存下标，按对应值保持单调。求左侧最近严格更小元素时，把值大于等于当前值的栈顶全部弹出；剩余栈顶就是答案，然后压入当前位置。被弹出的点比当前点更远，且值不更小，对后续位置不会比当前点更优，所以可以永久删除。

反向从右往左扫描，执行相同过程，就能求右侧最近严格更小位置。每个下标在每次扫描中至多入栈、出栈各一次，总时间为 $O(n)$；求更大位置则反转比较方向。
\subsubsection*{使用教程}
\InlineCode{init(a)} 预处理两侧最近严格更小的位置，\InlineCode{init(a,false)} 改为严格更大。结果为 \InlineCode{L[i],R[i]}；左侧不存在时为 0，右侧不存在时为 $n+1$。再次调用会覆盖上次结果。
\begin{lstlisting}[style=mtycpp,numbers=none]
vector<long long> a={0,3,1,1,4,2};
MonoStack st;
st.init(a);
int l=st.L[4],r=st.R[4];
st.init(a,false);
int next=st.R[2];
\end{lstlisting}
若要改为非严格关系，需要同时修改两遍扫描所使用的弹栈条件：
\[
\begin{array}{c|c}
\text{需要找到的关系}&\text{弹栈条件}\\ \hline
a_j<a_i&a_{top}\ge a_i\\
a_j\le a_i&a_{top}>a_i\\
a_j>a_i&a_{top}\le a_i\\
a_j\ge a_i&a_{top}<a_i
\end{array}
\]
\needspace{7\baselineskip}
\subsubsection*{模板代码}
\CodeFile{code/data-structures/monotonic_stack.cpp}
\subsubsection*{需要注意}
相等元素是否弹出决定严格与非严格，不能凭“递增栈”几个字判断条件。若用左右边界统计每个元素作为区间最小值的贡献，重复值通常需要一侧严格、另一侧非严格，避免同一区间被重复统计；本模板默认两侧都严格，不能直接把结果乘起来用于该贡献问题。笛卡尔树的栈构造见第~\ref{tree:cartesian}~节。
